\chapter{Архитектура заряженного посредника бинарной среды}
\label{ch:v8-charged-environment-mediator-architecture}

\section{Минимальная комплексная дилатация}

Пусть системный бинарный носитель имеет базис $|u\rangle,|d\rangle$,
гиперзаряд и гамильтониан
\begin{equation}
 Y_S=\operatorname{diag}(5/3,2/3),\qquad
 H_S=7\gamma |u\rangle\langle u|.
 \label{eq:v8-charged-mediator-system-data}
\end{equation}
Минимальная комплексная среда для downward-перехода есть
\begin{equation}
 \mathcal K_{\mathbb C}=\mathbb C|0\rangle\oplus\mathbb C|+\rangle,
 \quad Y_E=|+\rangle\langle+|,
 \quad H_E=7\gamma |+\rangle\langle+|.
 \label{eq:v8-charged-mediator-complex-environment}
\end{equation}
Для $b_+=|+\rangle\langle0|$ и $L_\downarrow=|d\rangle\langle u|$
зададим
\begin{equation}
 G=g\left(L_\downarrow\otimes b_+
 +L_\uparrow\otimes b_+^\dagger\right).
 \label{eq:v8-charged-mediator-interaction}
\end{equation}
Компенсация зарядов и резонанс дают два точных закона сохранения:
\begin{equation}
 [Y_S\otimes1+1\otimes Y_E,G]=0,\qquad
 [H_S\otimes1+1\otimes H_E,G]=0.
 \label{eq:v8-charged-mediator-conservation}
\end{equation}
Оператор $G/g$ имеет спектр $\{-1,0,0,1\}$ и действует нетривиально
только на резонансной плоскости
\begin{equation}
 \mathcal R=\operatorname{span}\{|u,0\rangle,|d,+\rangle\}.
 \label{eq:v8-charged-mediator-resonant-plane}
\end{equation}

\section{Точный канал одного столкновения}

Для $U_\theta=\exp[-i\theta G/g]$ и начального состояния $|0\rangle$
редуцированный канал имеет операторы Крауса
\begin{equation}
 K_0=P_d+\cos\theta\,P_u,\qquad
 K_1=-i\sin\theta\,L_\downarrow,\qquad
 K_0^\dagger K_0+K_1^\dagger K_1=1.
 \label{eq:v8-charged-mediator-kraus}
\end{equation}
Он является точным amplitude-damping каналом:
\begin{equation}
 \rho\longmapsto
 \begin{pmatrix}
 \cos^2\theta\,\rho_{uu}&\cos\theta\,\rho_{ud}\\
 \cos\theta\,\rho_{du}&\rho_{dd}+\sin^2\theta\,\rho_{uu}
 \end{pmatrix},
 \qquad p_\theta=\sin^2\theta.
 \label{eq:v8-charged-mediator-channel}
\end{equation}
При $0<p_\theta<1$ его ранг Крауса равен двум. Поэтому комплексная
размерность среды не может быть меньше двух, и конструкция минимальна.

\section{Конечная среда не создаёт необратимость}

Одно и то же возбуждение среды возвращает амплитуду системе. Вероятность
перехода периодична:
\begin{equation}
 p_{\theta+\pi}=p_\theta,\qquad
 p_{\pi/2}=1,\qquad p_\pi=0.
 \label{eq:v8-charged-mediator-revival}
\end{equation}
Это не полугруппа. Например, при $\theta=\pi/6$ композиция двух уже
редуцированных каналов имеет параметр
\begin{equation}
 p\star p=1-(1-p)^2=7/16,\qquad p=1/4,
 \label{eq:v8-charged-mediator-channel-composition}
\end{equation}
тогда как единая когерентная эволюция до $2\theta=\pi/3$ даёт
\begin{equation}
 p_{2\theta}=3/4,\qquad p_{2\theta}-(p\star p)=5/16.
 \label{eq:v8-charged-mediator-semigroup-defect}
\end{equation}
Необратимый предел требует свежей цепи вспомогательных систем и слабого
масштабирования
\begin{equation}
 \theta_h=g\sqrt h,\qquad
 \Phi_h=\operatorname{id}+h g^2\mathcal D_{L_\downarrow}+O(h^{3/2}).
 \label{eq:v8-charged-mediator-weak-collision}
\end{equation}
Ни цепь, ни механизм сброса, ни связь $h$ с физическими часами не следуют
из двухуровневой дилатации.

\section{Real-замыкание}

Пара зарядов $\{0,+1\}$ не замкнута относительно Real-сопряжения.
Минимальное Real-совместимое завершение имеет вид
\begin{equation}
 \mathcal K_R^{\mathbb C}=\mathbb C|0\rangle\oplus
 \mathbb C|+\rangle\oplus\mathbb C|-\rangle,\qquad
 Y_E^R=\operatorname{diag}(0,1,-1),\qquad J_E|+\rangle=|-\rangle.
 \label{eq:v8-charged-mediator-real-completion}
\end{equation}
Обе charged-линии должны быть семейными синглетами. Это новый тип данных:
старый трёхкратный charged carrier не содержит канонической неподвижной
линии.

\begin{remark}
Последующий admission-гейт уточняет эту coarse-классификацию: после
разложения multiplicity как $\mathbf1\oplus\mathbf3$ старая линия
$e_R^{(s)}$ и её Real-сопряжение уже дают charged singlet-пару. Новыми
остаются нейтральный carrier и динамические данные.
\end{remark}

Минимальный пакет новых данных равен
\begin{equation}
 \mathcal D_{\rm new}=\{\mathcal K_{+1}^{\rm sing}\oplus
 \mathcal K_{-1}^{\rm sing},\,7\gamma,\,|0\rangle,\,g,\,
 \text{fresh ancilla chain},\,\text{clock--rate anchor}\}.
 \label{eq:v8-charged-mediator-new-data}
\end{equation}
Итоговые реестры:
\begin{equation}
 N_{\rm complex\ dilation}=8/8,\qquad
 N_{\rm gauge/energy/Real}=6/6,\qquad
 N_{\rm irreversible\ origin}=0/6.
 \label{eq:v8-charged-mediator-ledgers}
\end{equation}

\begin{theorem}[Минимальный charged mediator]
Двухуровневая комплексная среда с зарядами $0,+1$ и щелью $7\gamma$
даёт минимальную энергосохраняющую и gauge-инвариантную дилатацию точного
amplitude-damping канала за одно столкновение. Real-замыкание требует
третьей линии заряда $-1$. Конечная дилатация остаётся периодической и не
порождает необратимую полугруппу без свежей цепи и слабого предела.
\end{theorem}

\begin{nogocriterion}[Дилатация одного шага не есть динамический parent]
Нельзя отождествлять точный канал одного столкновения с автономной
необратимой динамикой. Нельзя удалять Real-сопряжённую charged-линию и
нельзя считать family-singlet назначение унаследованным из старого
триплета. Параметры $g$, состояние среды и clock--rate anchor остаются
новыми входами.
\end{nogocriterion}

\begin{protocol}\begin{enumerate}
\item минимальная комплексная размерность среды: \textbf{$2$};
\item заряды комплексной среды: \textbf{$0,+1$};
\item резонансная щель: \textbf{$7\gamma$};
\item полный заряд сохраняется: \textbf{да};
\item свободная энергия сохраняется: \textbf{да};
\item канал одного столкновения: \textbf{точный amplitude damping};
\item ранг Крауса внутри $0<p<1$: \textbf{$2$};
\item минимальная Real-совместимая размерность: \textbf{$3$};
\item одно конечное окружение задаёт полугруппу: \textbf{нет};
\item точный композиционный дефект при $\theta=\pi/6$: \textbf{$5/16$};
\item complex dilation: \textbf{$8/8$};
\item gauge/energy/Real completion: \textbf{$6/6$};
\item autonomous irreversible origin: \textbf{$0/6$};
\item следующий гейт: \textbf{admission существующего charged carrier}.
\end{enumerate}\end{protocol}

Машинный сертификат сохранён в
\path{s2t_v8_baryon_c0_singlet_triplet_central_gap_isotypic_channel_extra_edge_mass_two_source_cubic_incidence_up_down_binary_selector_minimal_typed_discriminator_binary_carrier_charged_environment_mediator_architecture_gate_results.json}.

\begin{thebibliography}{9}
\bibitem{v8-charged-mediator-attal} S.~Attal and Y.~Pautrat,
\emph{From Repeated to Continuous Quantum Interactions},
Ann. Henri Poincar\'e 7 (2006) 59; arXiv:math-ph/0311002.
\bibitem{v8-charged-mediator-haapasalo} E.~Haapasalo,
\emph{The Choi--Jamio\l{}kowski Isomorphism and Covariant Quantum Channels},
Quantum Stud. Math. Found. 8 (2021) 1; arXiv:1906.11442.
\end{thebibliography}